%!TEX root = ../../../exa-ma-deliverables.tex
\VThreeApplicationDraft{app-feelpp-mi300a-fem}
{MI300A high-order finite-element kernels and HiFiMagnet}
{C. Prud'homme; V. Chabannes; P. Jolivet}
{WP1, WP3, WP7}{extended-mini-app}
{Benchmark the forward nonlinear thermoelectric high-field-magnet model at fixed accuracy, including solver behaviour, device residency and CPU/APU performance.}
{Feel++, PETSc, HPDDM; MPI, Kokkos and HIP}{implementation-in-progress}
\subsection{Description of the benchmark}
Forward-performance benchmark: solve the nonlinear thermo-electric high-field-magnet model at fixed accuracy and expose state fields, nonlinear/linear solver behaviour, device residency and CPU/APU performance; no inverse or control objective is introduced in this application.
\paragraph{Scope boundary.}
This record covers only the shared forward operator.  Parameter inversion and
optimal control are separate applications.  The Sheet describes a future MI300A
campaign; this integration draft does not present that allocation or campaign as
a completed result.


\subsection{Benchmarking methodology}


\paragraph{Evidence and benchmark contract.}
The inspected \href{https://github.com/numpex/apps-feelpp/commit/bf1febf41cd615c20928e1bb4ef70678533c0587}{application commit}
contains a HiFiMagnet thermoelectric JSON configuration under
\texttt{src/app-feelpp-toolbox/thermoelectric/hifimagnet}.  It does not yet
establish a dedicated executable benchmark or MI300A result for this ID.
\texttt{BM-FEEL-MI300A-01} therefore remains implementation in progress.
\todo{Pin the executable/config, input data, native and container environments,
MI300A topology, fixed-accuracy baseline, repetitions and raw-data location.}
\begin{description}
  \item[Method-Algorithm WP1:] high-order cG/H1 and HDG, full/partial/matrix-free assembly, homogeneous worksets, facets, traces, coupled blocks, skeleton terms
  \item[Method-Algorithm WP3:] PETSc/HPDDM Krylov solvers, preconditioning, nonlinear and block solvers, device-resident sparse algebra
  \item[Framework:] Feel++, PETSc, HPDDM
  \item[parallel\_framework:] MPI, Kokkos, GPU - HIP, Multithread - OpenMP
  \item[Benchmark scope:] single-node, multi-node, multi-gpu, full-system, method-verification, solver-scaling
  \item[metrics:] forward-residual, discretization-error, fixed-accuracy, time-to-solution, time-to-precision, nonlinear-and-linear-iterations, assembly-cost, solve-cost, memory-residency, host-device-transfers, strong-scalability, weak-scalability, multi-APU-scalability
\end{description}
\todo[app-feelpp-mi300a-fem / benchmarking methodology]{Benchmarking methodology to be supplied by contributors.}

\subsection{Input/Output Dataset Description}

\subsubsection{Input Data}
Gmsh, JSON config, Spack environment, Apptainer image

\subsubsection{Output Data}
Electric potential $\varphi$, temperature T, unknown potential drop $\Delta$V, current density J, Joule losses, JSON/CSV performance reports, profiles, HDF5 fields, reproducible native and container environments

\subsection{Results summary}
\todo[app-feelpp-mi300a-fem / results summary]{Results summary to be supplied by contributors.}
